\section{Translating into Lean}\label{sec:10-modules:02-translating-into-lean:1}

The same ``data, then axioms'' build used for \texttt{Group} and \texttt{Ring} applies here:

\begin{lstlisting}
structure Module (R : Type) (Rg : Ring R) (M : Type) where
  addGrp : CommGroup M
  smul : R → M → M
  smul_add : ∀ (r : R) (m n : M), smul r (addGrp.op m n) = addGrp.op (smul r m) (smul r n)
  add_smul : ∀ (r s : R) (m : M), smul (Rg.addGrp.op r s) m = addGrp.op (smul r m) (smul s m)
  smul_smul : ∀ (r s : R) (m : M), smul (Rg.mul r s) m = smul r (smul s m)
  one_smul : ∀ m : M, smul Rg.one m = m
\end{lstlisting}

Field by field:

\begin{itemize}
\tightlist
\item
  \texttt{addGrp : CommGroup M} --- the underlying abelian group of the module, exactly as \texttt{addGrp} played this role inside \texttt{Ring} (Chapter 8). Note that \texttt{Module} takes \texttt{Rg : Ring R} as an \emph{explicit argument}, not a field: a module is always a module \emph{over} a specific, already-given ring, so \texttt{Rg} is a parameter of the whole structure rather than data bundled inside it.
\item
  \texttt{smul : R → M → M} --- the scalar action, \texttt{r • m} in ordinary notation.
\item
  \texttt{smul\_\allowbreak{}add}, \texttt{add\_\allowbreak{}smul} --- the two distributivity laws (M1), (M2). Read them as ``scalar over module-sum'' and ``ring-sum over scalar''. This is the same left/right split used for \texttt{Ring}'s \texttt{left\_\allowbreak{}distrib}/\texttt{right\_\allowbreak{}distrib}, but here the split is based on \emph{which} addition (\texttt{Rg.addGrp.op} vs. \texttt{addGrp.op}) is involved, not on which side \texttt{mul} is applied.
\item
  \texttt{smul\_\allowbreak{}smul} --- (M3), compatibility of the ring's multiplication with iterated scalar action.
\item
  \texttt{one\_\allowbreak{}smul} --- (M4), the ring's multiplicative identity acts as the identity scalar.
\end{itemize}

\begin{mathreading}

\texttt{Module R Rg M} is the type of left \(R\)-module structures on the abelian group \(M\). The data is an action \(R \times M \to
M\), \((r,m)\mapsto r\cdot m\). The four axioms say precisely that the curried map \(\rho(r)(m) = r\cdot m\) is a \textbf{ring homomorphism} from \(R\) into the ring of endomorphisms of \((M,+)\): the additive maps \(M \to M\), under pointwise addition and composition. \texttt{smul\_\allowbreak{}add} says each \(\rho(r)\) is itself additive (a homomorphism \(M \to M\) of abelian groups). \texttt{add\_\allowbreak{}smul} and \texttt{smul\_\allowbreak{}smul} say \(\rho\) preserves \(+\) and \(\times\). And \texttt{one\_\allowbreak{}smul} says \(\rho(1)\) is the identity map on \(M\). So a module over \(R\) is exactly an abelian group \(M\) equipped with a ring homomorphism from \(R\) into its own ring of endomorphisms. The ring \(R\) enters as an explicit \emph{parameter}, not as extra bundled data on \(M\), because a module is always a module \emph{over} some already-fixed ring.

\end{mathreading}

\begin{quote}
Read more: Mathlib's \href{https://loogle.lean-lang.org/?q=Module}{\texttt{Module}} (\texttt{Mathlib.Algebra.Module.Defs}) is much more general. It is universe-polymorphic, stated for \texttt{Semiring} rather than just \texttt{Ring}, and integrated with the whole \texttt{Mathlib.LinearAlgebra.*} hierarchy (bases, dimension, tensor products); see \hyperref[sec:13-next-steps:02-moving-to-mathlib:1]{Chapter 13}. For the classical theory, any standard module-theory or homological-algebra text (e.g.~Dummit \& Foote's chapters on modules, or Weibel's \emph{An Introduction to Homological Algebra} for the deeper theory) covers the same axioms from the ground up.
\end{quote}

\subsection{References}\label{sec:10-modules:02-translating-into-lean:2}

Full citations in the \hyperref[sec:root:bibliography:1]{Bibliography}.

\begin{itemize}
\tightlist
\item
  Dummit and Foote (\cite{DummitFoote2003}) --- the module axioms (M1)--(M4) from the classical, non-Lean point of view.
\item
  Weibel (\cite{Weibel1994}) --- the deeper module/homological theory this section's ``Read more'' points toward.
\end{itemize}
